\fsection{Tipo de instalación y ubicación}
Se trata de una celda robotizada de manipulación (pick and place) para logística y preparación de pedidos,
dentro del sector del comercio electrónico. El prototipo se monta y valida en el aula-taller de Automatización y
Robótica Industrial del IES Jorge Manrique, en Motilla del Palancar (Cuenca), sobre el equipamiento Robot Pick
AI del que dispone el centro gracias al proyecto de innovación en el que participa junto con LH FP Barakaldo,
el IES Javier García Téllez de Cáceres y la empresa Fluitronic S.L.

El entorno es una instalación interior, a temperatura ambiente y sin requisitos especiales de protección
ambiental. Todos los equipos se comunican por una red Ethernet industrial (Profinet) propia de la celda, en el
rango 192.168.15.xxx y sin conexión a servicios en la nube: el modelo de IA se ejecuta en local en el IPC. Al
usarse un robot colaborativo, la celda puede compartir espacio con los operarios, aunque se realiza la
correspondiente evaluación de riesgos de la aplicación.
